The baselines set three kinds of reference for the generative model.

\paragraph{What is already easy.} MINOS matching is almost a deterministic function of the true muon
kinematics: the tree reaches AUC~0.987 and cuts the log loss from 0.52 to 0.10. The reconstruction
efficiency is predicted without visible bias across muon angle, muon momentum and hadron count
(Fig.~\ref{fig:eff}); its residual log loss (0.42 versus 0.67 marginal) is irreducible randomness in
whether a muon track is found, not a modelling deficit. The recoil-energy median follows the true hadronic
kinetic energy over two decades and the MDN reproduces the 16--84\% band (Fig.~\ref{fig:recoil}).

\paragraph{What needs a proper generative model.} The muon momentum response is bimodal in the sense that
matters for analyses: MINOS-matched muons have a $\pm 9\%$ core, unmatched muons a spread of more than a
factor two in either direction. The $K=1$ Gaussian cannot represent this (Table~\ref{tab:nll}); the
$K=8$ MDN does, with $W_1 = 0.008$ in $\log(P_\mathrm{reco}/P_\mathrm{true})$ for matched and $0.05$ for
unmatched muons. But the MDN sees only 48 summary features. Correlations between the muon response, the
recoil energy and the prong content of the same event are not modelled at all here, and these are exactly
what a cross-section analysis with a multi-variable signal definition depends on. That joint structure is the
job of the M2 flow-matching model over the particle set.

\paragraph{Where the baseline is weak.} Multiplicity is the hardest target so far. The tree beats the
marginal (log loss 0.84 versus 1.06, accuracy 0.62 versus 0.44) and reproduces the mean reco prongs versus
true charged hadrons up to three hadrons (Fig.~\ref{fig:mult}), but the confusion table shows that events
with four or more true charged hadrons are still mostly assigned one or two prongs, with the saturation at
high multiplicity slightly overshot. Whether the saturation is learnable from the particle set rather than
from counts is one of the questions the M2 model must answer; it also bears directly on the extrapolation
study, since high-multiplicity final states are where alternative generators differ most.

\paragraph{Caveats.} All numbers are from a single MC file (368k CC~$\nu_\mu$ events, 39.6\% reconstructed)
and a single seed. Test negative log-likelihoods for the muon MDN are quoted with the same $\pm 25\sigma$
robust clipping used in training; without it a handful of muons with $|p_x/p_z| \gg 1$ dominate the mean.
The features exclude generator-internal labels by design, so nothing here would change for another
generator's final states, but neither has that transfer been tested yet: the extrapolation holdouts of
the design document (train on QE+RES+2p2h, test on DIS; train on $\le 2$ visible hadrons, test on more)
are scheduled for M3.
